{\scriptsize
\begin{longtable}{p{0.45cm}p{1.85cm}p{0.95cm}p{2.0cm}p{2.0cm}p{2.0cm}p{2.0cm}p{2.0cm}}
\toprule
Ind. & Short name & Dim. & Score 0 & Score 1 & Score 2 & Score 3 & Interpretation \\
\midrule
\endfirsthead
\toprule
Ind. & Short name & Dim. & Score 0 & Score 1 & Score 2 & Score 3 & Interpretation \\
\midrule
\endhead
A & Export compensation friction & Revenue & Stable and attractive compensation (e.g., generous FiT or close-to-retail offset). & Moderate discount relative to retail value. & Clearly weaker export value than self-consumption, but still usable. & Very weak export compensation or highly unfavorable export pricing. & Higher values indicate weaker monetization of exported generation. \\
B & Export constraint friction & Revenue & No meaningful constraint on offset, carry-forward, or monetization. & Minor or occasional constraints. & Noticeable limits on export use or contractual caps. & Strong limits such as no cash-out, no carry-forward, or strict export-capacity ceilings. & Captures whether exported electricity can be converted into durable economic value. \\
C & Settlement complexity friction & Revenue & Simple and predictable settlement (e.g., straightforward FiT or netting). & Mostly standardized settlement with limited complexity. & Multiple contractual or tariff steps materially affect settlement. & Highly complex settlement involving market exposure, layered contracts, or difficult bill reconciliation. & Captures whether monetization depends on complicated billing or market arrangements. \\
D & Policy uncertainty friction & Revenue & Stable long-term policy with clear continuation horizon. & Mostly stable framework with limited uncertainty. & Meaningful uncertainty due to reform risk or short policy horizons. & High uncertainty due to active reform, marketization, or unstable support mechanisms. & Captures how predictable future rooftop PV revenue streams are. \\
E & Small-system approval friction & Administrative & Clear simplified path or self-issue route for small systems. & Mostly simple approval process, with some documentation requirements. & Formal approval remains necessary and simplification is limited. & No meaningful simplification; small systems face full approval burden. & Designed to distinguish jurisdictions with streamlined residential permitting from those without. \\
F & Building/planning approval friction & Administrative & Minimal additional building/planning review in ordinary cases. & Some additional review only for specific conditions (e.g., heritage, structural change). & Routine administrative review or filing is typically required. & Multi-agency or substantial building/planning review is commonly required. & Captures urban-form and building-governance barriers beyond pure energy regulation. \\
G & Grid study/fee friction & Administrative & No material fee or study burden for typical rooftop systems. & Light utility processing or limited technical checks. & Meaningful fees, studies, or multiple utility steps for many systems. & Heavy study requirements, high fees, or substantial interconnection burden. & Reflects utility-facing transaction costs and soft costs. \\
H & Professional credential friction & Administrative & No unusual credential requirement beyond normal electrical practice. & Some certified participation is advisable or needed in limited cases. & Certified or licensed professionals are routinely required. & Multiple formal credentials, recognized intermediaries, or tightly restricted installer channels are required. & Captures labor-market and procedural gatekeeping. \\
\bottomrule
\end{longtable}
}
